% !TeX spellcheck = ru_RU
\documentclass[a4paper, 14pt, unknownkeysallowed]{extreport}
\usepackage[utf8]{inputenc}
\usepackage[T2A]{fontenc}
\usepackage[english,russian]{babel}
\usepackage{cmap}
\usepackage{enumitem}

\usepackage{csvsimple}
\usepackage{hyphenat} 
\usepackage{amstext, amsmath,amsfonts,amssymb,amsthm,mathtools} 

\usepackage{setspace}
\onehalfspacing 

\usepackage{geometry}
\geometry{left=30mm}
\geometry{right=10mm}
\geometry{top=10mm}
\geometry{bottom=20mm}

\usepackage{graphicx}
\graphicspath{{img/}} 

\usepackage{threeparttable}
\usepackage{bigdelim}

\usepackage{setspace}
\onehalfspacing % Полуторный интервал

\frenchspacing
\usepackage{indentfirst} % Красная строка

\usepackage{cases}

\usepackage[unicode,pdftex]{hyperref} % Ссылки в pdf
\hypersetup{hidelinks}

\usepackage{xcolor}
\usepackage{listings}
% Для листинга кода:
\lstset{%
	language=lisp,   					% выбор языка для подсветки	
	basicstyle=\small\sffamily,			% размер и начертание шрифта для подсветки кода
	numbers=left,						% где поставить нумерацию строк (слева\справа)
	%numberstyle=,					% размер шрифта для номеров строк
	stepnumber=1,						% размер шага между двумя номерами строк
	numbersep=5pt,						% как далеко отстоят номера строк от подсвечиваемого кода
	frame=single,						% рисовать рамку вокруг кода
	tabsize=4,							% размер табуляции по умолчанию равен 4 пробелам
	captionpos=t,						% позиция заголовка вверху [t] или внизу [b]
	breaklines=true,					
	breakatwhitespace=true,				% переносить строки только если есть пробел
	escapeinside={\#*}{*)},				% если нужно добавить комментарии в коде
	backgroundcolor=\color{white},
}
\usepackage{caption}
\captionsetup{labelsep = endash}
\captionsetup[figure]{name = {}, justification = centerlast,labelformat=empty}
\captionsetup[lstlisting]{justification = justified}
\captionsetup[table]{justification = justified}


\usepackage{titlesec}
\titleformat{\chapter}{\LARGE\bfseries}{\thechapter}{20pt}{\LARGE\bfseries}
\titleformat{\section}{\Large\bfseries}{\thesection}{20pt}{\Large\bfseries}

\renewcommand\labelitemi{\ --\ }
\renewcommand\labelenumi{\theenumi)}

\newcommand{\img}[3] {
	\begin{figure}[h!]
		\includegraphics[scale = #1]{img/#2}
		\caption{#3}
		\label{img:#2}
	\end{figure}
}

\newcommand{\lst}[4]{
\lstinputlisting[language=Lisp, firstline=#1, lastline=#2, label=lst:#3,caption=#4]{sx.lsp}
}

\usepackage{siunitx,array,booktabs}


\begin{document}
% !TeX spellcheck = ru_RU
\begin{titlepage}
	
	\newgeometry{pdftex, left=2cm, right=2cm, top=2.5cm, bottom=2.5cm}
	\fontsize{12pt}{12pt}\selectfont
	\noindent \begin{minipage}{0.15\textwidth}
		\includegraphics[width=\linewidth]{img/b_logo.jpg}
	\end{minipage}
	\noindent\begin{minipage}{0.9\textwidth}\centering
		\textbf{Министерство науки и высшего образования Российской Федерации}\\
		\textbf{Федеральное государственное бюджетное образовательное учреждение высшего образования}\\
		\textbf{«Московский государственный технический университет имени Н.~Э.~Баумана}\\
		\textbf{(национальный исследовательский университет)»}\\
		\textbf{(МГТУ им. Н.~Э.~Баумана)}
	\end{minipage}

	\noindent\rule{18cm}{3pt}
	\newline\newline
	\noindent ФАКУЛЬТЕТ $\underline{\text{«Информатика и системы управления»~~~~~~~~~~~~~~~~~~~~~~~~~~~~~~~~~~~~~~~~~~~~~~~~~~~~~~~}}$ \newline\newline
	\noindent КАФЕДРА $\underline{\text{«Программное обеспечение ЭВМ и информационные технологии»~~~~~~~~~~~~~~~~~~~~~~~}}$\newline\newline\newline\newline\newline\newline\newline
	
	
	\begin{center}
		\Large\textbf{Исследование быстродействия умножения}
		
		\Large\textbf{матриц на языке Common Lisp}
%		
		\textbf{\newline}		
	\end{center}
	
	\noindent\textbf{Студент} $\underline{\text{~~Виноградов А. О.~~~~~~~~~~~~~~~~~~~~~~~~~~~~~~~~~~~~~~~~~~~~~~~~~~~~~~~~~~~~~~~~~~~~~~~~~~~~~~~~~~~~~~}}$\newline\newline
	\noindent\textbf{Группа} $\underline{\text{~~ИУ7-66Б~~~~~~~~~~~~~~~~~~~~~~~~~~~~~~~~~~~~~~~~~~~~~~~~~~~~~~~~~~~~~~~~~~~~~~~~~~~~~~~~~~~~~~~~~~~~~~~~~~~~}}$\newline\newline
	\noindent\textbf{Преподаватель} $\underline{\text{~~Строганов Ю.В.~~~~~~~~~~~~~~~~~~~~~~~~~~~~~~~~~~~~~~~~~~~~~~~~~~~~~~~~~~~~~~~~~~~~~~~~~~~~~}}$\newline
	
	\begin{center}
		\vfill
		Москва~---~\the\year
		~г.
	\end{center}
	\restoregeometry
	
\end{titlepage}
\setcounter{page}{2}

\chapter*{Введение}
\addcontentsline{toc}{chapter}{Введение}
Цель работы: приобрести навыки создания и использования функций пользователя в Lisp.

Задачи работы:
\begin{itemize}
	\item изучить работу интерпретатора Lisp;
	\item изучить алгоритм работы функции eval;
	\item изучить структуру и порядок обработки программы в Lisp.
\end{itemize}

\chapter{Практические задания}

\section{Составить диаграммы вычисления для \\выражений}

\img{1}{1_1}{}
\img{1}{1_2}{}
\img{1}{1_3}{}
\img{1}{1_4}{}
\img{1}{1_5}{}
\img{1}{1_6}{}

\clearpage
\section{Написать функцию, вычисляющую \\ гипотенузу прямоугольного треугольника по заданным катетам и составить диаграмму её вычисления}

\begin{lstlisting}
(defun hyp (a, b)
	(sqrt (+ (* a a) (* b b)))
)
\end{lstlisting}

\img{1}{2_1}{}

\section{Каковы результаты вычисления следующих выражений?}

\begin{lstlisting}
(list 'a c) ; unbound variable c, solution: (list 'a 'c)
(cons 'a (b c)) ; unbound variable (b c), solution: (cons 'a '(b c))
(caddr (1 2 3 4 5)) ; illegal function call, solution: (caddr '(1 2 3 4 5))
(cons 'a 'b 'c) ; invalid number of arguments: 3, solution:  (cons 'a (cons'b 'c))
(list 'a (b c)) ; unbound variable b, solution: (list 'a '(b c))
(list a '(b c)) ; unbound variable a, solution: (list 'a '(b c))
(list (+ 1 '(length '(1 2 3)))) ; type error, solution: (list (+ 1 (length '(1 2 3))))
\end{lstlisting}

\section{Написать функцию longer\_then от двух\\ списков-аргументов, которая возвращает Т, если первый аргумент имеет большую\\ длину}

\begin{lstlisting}
(defun longer_then (a b)
	(> (length a) (length b))
\end{lstlisting}

\section{Каковы результаты вычисления следующих выражений?}

\begin{lstlisting}
(cons 3 (list 5 6))	;(3 5 6)
(cons 3 '(list 5 6)) ;(3 list 5 6)
(list 3 'from 9 'lives (- 9 3))	;(3 from 9 lives 6)
(+ (length for 2 too)) (car '(21 22 23))) ;the variable FOR is unbound
(cdr '(cons is short for ans))	;(is short for ans)
(car (list one two))	;the variable ONE is unbound
(car (list 'one 'two)) ;one
\end{lstlisting}

\section{Какие результаты вычисления следующих выражений?}
Дана функция (defun mystery (x) (list (second x) (first x))).

\begin{lstlisting}
(mystery (one two)) ;undefined function ONE
(mystery (last one two)) ;undefined function LAST
(mystery free) ;the variable FREE is unbound
(mystery one 'two) ;the variable ONE is unbound
\end{lstlisting} 

\section{Написать функцию, которая переводит \\температуру в системе Фаренгейта в \\температуру по Цельсию.}

\begin{lstlisting}
(defun f-to-c (x) 
	(* (/ 5 9) (- x 32)))
\end{lstlisting} 

Как бы назывался роман Р.Брэдбери "+451 по Фаренгейту" в системе по Цельсию?
Ответ: 232.78 по Цельсию

\section{Что получится при вычисления каждого из выражений?}

\begin{lstlisting}
(list 'cons t NIL)	;(cons t Nil)
(eval (list 'cons t NIL)) ;(T)
(eval (eval (list 'cons t NIL))) ;function T is undefined
(apply #'cons '(t NIL))	;(T)
(eval NIL) ;Nil
(list 'eval NIL) ;(eval Nil)	
(eval (list 'eval NIL)) ;(Nil)
\end{lstlisting} 

\end{document}